\documentclass[11pt,a4paper,spanish]{book}
\usepackage{estilo_unir-1}
\usepackage{apacite}

\usepackage{hyperref}
\numberwithin{equation}{chapter}
%\usepackage{caption}
\numberwithin{figure}{chapter}
%\usepackage{chngcntr}
%\counterwithin{equation}{chapter}
%\counterwithin{figure}{chapter}
%\counterwithin{table}{chapter}
%\renewcommand\theequation{\thechapter.\arabic{equation}}
%\renewcommand\thefigure{\thechapter.\arabic{figure}}
%\renewcommand\thetable{\thechapter.\arabic{table}}
%\renewcommand\theequation{\thechapter.\arabic{equation}}
%\counterwithin{figure}{chapter}
%\counterwithin{table}{chapter}
%\renewcommand\thefigure{\thechapter.\arabic{figure}}
%\renewcommand\thetable{\thechapter.\arabic{table}}
%\renewcommand\thefigure{\thechapter.\arabic{figure}}
%\makeatletter
%\renewcommand\p@figure{\thechapter..\arabic{figure}}
%\makeatother



%---------------------------
%título del trabajo y autor
%---------------------------
\title{Reconocimiento de patrones en sistemas industriales para la detección de fallos
		mediante algoritmos de Aprendizaje Automático}
\titulacion{Máster Universitario en Inteligencia Artificial}
\author{Jorge Elicer Lambraño Arroyo}
\date{11 de Septiembre de 2025}
\director{ }
\nombreciudad{Bogotá (Colombia)}

%---------------------------
%marges
%---------------------------
%\usepackage[margin=1.9cm]{geometry}
%---------------------------
%---------------------------
%---------------------------
%---------------------------
\begin{document}
\renewcommand{\listfigurename}{Índice de Ilustraciones}
\renewcommand{\listtablename}{Índice de Tablas}
\renewcommand{\contentsname}{Índice de Contenidos}
\renewcommand{\figurename}{Figura}
\renewcommand{\tablename}{Tabla} 

\maketitle

\frontmatter
\tableofcontents
\listoffigures
\listoftables

\chapter{Resumen}
% {\bf Nota:} En este apartado se introducirá un breve resumen en español del trabajo 
% realizado (extensión máxima: 150 palabras). Este resumen debe incluir el objetivo o 
% propósito de la investigación, la metodología, los resultados y las conclusiones.
El presente proyecto aborda la temática de  la detección de anomalías en rodamientos 
utilizando el CWRU Bearing Dataset como fuente de información. El objetivo principal 
consiste en comparar el desempeño de diversos modelos de aprendizaje automático con el 
fin de identificar cuál resulta más apropiado para el desarrollo de un detector de 
anomalías robusto, eficiente y aplicable en entornos industriales. Para ello, se 
implementaron diferentes etapas de preprocesamiento,  caracterización de las señales 
temporales y estrategías durante la fase el entrenamiento de los modelos, con el 
propósito de obtener un modelo optimizado con altas métricas de rendimiento que facilite
la tarea de detección de anomalías.Los resultados alcanzaron métricas superiores al 
90\% de exactitud, precisión, sensibilidad, F1-score y área bajo la curva RoC, lo que 
demuestra la viabilidad de aplicar machine learning en la detección temprana de fallas
en rodamientos

{\bf Palabras Clave:} Aprendizaje automático; CWRU Bearing Dataset; Detección de fallas;
Mantenimiento predictivo; Preprocesamiento de datos.

\chapter{Abstract}
% {\bf Nota:} En este apartado se introducirá un breve resumen en español del trabajo
% realizado (extensión máxima: 150 palabras). Este resumen debe incluir el objetivo o 
% propósito de la investigación, la metodología, los resultados y las conclusiones.
The present project addresses the topic of anomaly detection in bearings using the CWRU
Bearing Dataset as the source of information. The main objective is to compare the 
performance of various machine learning models in order to identify which one is the most
suitable for the development of a robust, efficient, and industrially applicable anomaly
detector. To this end, different preprocessing stages, time-series signal characterization,
and training strategies were implemented with the purpose of obtaining an optimized model
with high performance metrics that facilitates the anomaly detection task. The results 
achieved over 90\% in accuracy, precision, recall, F1-score and RoC curve, demonstrating
the feasibility of applying machine learning for early fault detection in bearings.

{\bf Palabras Clave:} Machine Learning; CWRU Bearing Dataset; Fault detectino;
Predictive maintenance; Data preprocessing.




\mainmatter
\chapter{Introducción}

% El primer capítulo es siempre una introducción. En la introducción se debe resumir 
% de forma esquemática % pero suficientemente clara lo esencial de cada una de las partes 
% del trabajo. La lectura de este primer capítulo ha de dar una idea clara de lo que se 
% pretendía, las conclusiones a las que se ha llegado y del procedimiento seguido.

% Como tal, es uno de los capítulos más importantes de la memoria. Las ideas principales
% a transmitir son la identificación del problema a tratar, la justificación de su 
% importancia, los objetivos generales a grandes rasgos y un adelanto de la contribución
% que esperas hacer.

% Típicamente una introducción tiene tres apartados:
% \begin{itemize}
% \item Motivación / justificación del tema a tratar
% \item Planteamiento del trabajo
% \item Estructura del trabajo
% \end{itemize}






\chapter{Contexto y Estado del Arte}

\chapter{Identificación de Requisitos}

\chapter{Objetivos}

\chapter{Desarrollo del trabajo}

\chapter{Conclusiones y Trabajo Futuro}

En el presente  trabajo se alcanzaron métricas superiores al 90\% en todos los modelos
evaluados, superando el resultado esperado, lo cual refleja la solidez del enfoque 
planteado para la detección de anomalías en entornos industriales. La etapa de 
preprocesamiento de los datos fue un paso clave: la caracterización de las series de 
tiempo, junto con un adecuado balanceo del dataset y la búsqueda de hiperparámetros, 
permitió que los modelos pudieran generalizar mejor y mostraran robustez frente al 
sobreajuste. Estos resultados confirman la efectividad del pipeline desarrollado para 
transformar datos crudos en características discriminativas útiles para los modelos 
de predicción.


Debido al alto rendimiento alcanzado, todos los modelos evaluados podrían aplicarse en 
una solución in-situ. No obstante, los que obtuvieron las mejores métricas fueron 
XGBoost y LightGBM, confirmando su eficacia como algoritmos de referencia en problemas
de clasificación complejos, y un buen referente para la detección de anomalías. Sin 
embargo, si la implementación requiere simplicidad, interpretabilidad y bajo costo 
computacional, los modelos clásicos como la regresión logística, los árboles de decisión
o Naive Bayes resultan altamente recomendables, ya que ofrecen un equilibrio entre 
eficiencia, precisión e interpretabilidad.


Finalmente, gracias al uso de modelos con mayor interpretabilidad, se identificó que 
variables como los valores pico (maximum), la desviación estándar (std), el valor RMS,
la frecuencia fundamental (f0) y la distorsión armónica total (THD) son las que más 
aportan a la clasificación de una observación como anomalía. Este hallazgo  además de 
validar la calidad del modelo, proporciona información valiosa para la toma de 
decisiones en el ámbito industrial, al permitir inferir cuáles son los factores físicos
más determinantes en el diagnóstico temprano de fallos.


Como línea de trabajo futuro, una posibilidad interesante es la integración de los 
modelos desarrollados en un PLC (Controlador Lógico Programable) u otro dispositivo
de cómputo industrial. Esto permitiría que la detección de anomalías se ejecute de 
manera automática y en tiempo real dentro del propio entorno productivo. La adaptación
de los algoritmos a este tipo de hardware representa un reto, especialmente en lo 
relacionado con la optimización del uso de recursos computacionales, pero a la vez abre
la puerta a aplicaciones industriales de gran impacto.


Otra línea de investigación a considerar es el diseño de un algoritmo de clasificación
que no solo detecte anomalías, sino que también identifique el tipo específico de falla:
un clasificador multiclase.  Este enfoque aportaría un valor agregado, ya que permitiría
pasar de un sistema de simple diagnóstico (anómalo/no anómalo) a un sistema de prognosis
y clasificación de fallas, capaz de diferenciar entre problemas la gravedad del fallo y 
el lugar donde se presenta. Con ello, las soluciones podrían evolucionar hacia un 
mantenimiento predictivo más preciso y especializado.


En la Ecuación \eqref{eq:eq1secCTF}


\begin{equation}\label{eq:eq1secCTF}
M=\begin{pmatrix}
	m_{11}&m_{12}\\
	m_{21}&m_{22}
\end{pmatrix}
\end{equation}

En la siguiente Tabla \ref{tab:tab1secCTF}

\begin{table}[h]
\centering
\begin{tabular}{|c|c|}
	\hline
	1 & 2 \\
	\hline
	22 & 11 \\
	\hline
\end{tabular}
\caption{Tabla 1}
\label{tab:tab1secCTF}
\end{table}

En la siguiente Figura \ref{fig:fig1secCTF}

\begin{figure}[h]
\includegraphics[width= 0.8\textwidth]{logo_unir}
\caption{Logo Unir}
\label{fig:fig1secCTF}
\end{figure}

\cite{PIMENTEL2016744} \cite{da_S_Bessa_2023}

%\begin{thebibliography}{a}
%\bibitem{etiqueta} \textsc{Autores},
%\textit{nombre referencia.}
%Información addicional
%\end{thebibliography}
\bibliographystyle{apacite}
\bibliography{bibliografia}

\appendix
\chapter{Apendices}

\end{document}





















